\begin{resumo}[Abstract]
 \begin{otherlanguage*}{english}
  
 Artificial cellular organisms are circuits able to encapsulate and control integrative architectures, capable of performing all required calculations on the chosen discrete domain. Cellular automata are an alternative for modeling physical systems in a discrete domain. In this work, the physical phenomenon considered, as a case study, is the sound wave propagation through an ideal and stationary medium. The work was carried out in two stages, the first focused on the development of a CA system, applied to one-dimensional plane wave propagation, in a reconfigurable circuitry, and the second focused on its on-chip implementation proposals. The developed CA system, called Camphs1D, is composed of two automata cells, based on IP-cores of 27-bit arithmetic operators, capable of operate all basic wave phenomena in one-dimensional media. The aforementioned cells were developed from two different strategies, the first one directly described in VHDL, using Vivado software, and the second one using C/C++, using Vivado HLS software, where Xilinx 32-bit floating point IP-cores were used. In order to validate results, simulations were used in Testbench VHDL, simulations in SDK software (Xilinxs), and physical implementations integrated with a software ARM terminator and DDR memory access. All validations achieved their goals, when compared with the reference models, developed from the solution of d'Alembert for the one-dimensional wave equation.
 
 \vspace{\onelineskip}
 
 \noindent 
 \textbf{Key-words}: cellular automaton; digital waveguide mesh; systems on chip.
 \end{otherlanguage*}
\end{resumo}
